\section{Primordial Solar Monotheism}

This is the concluding section of the Arctic Myth chapter of Julius Evola's The Myth of Blood. Evola extracts two main points from Wirth:

\begin{enumerate}
\item The Arctic Myth
\item Nordic primordial monotheism
\end{enumerate}

The former Evola accepts, though not all of the latter. Nevertheless, today the Arctic Myth is more difficult to defend and a primordial monotheistic tradition seems plausible. For the Arctic myth, we can rely on Fabre d'Olivet, Bal Tilak, Hermann Wirth, Guenon and Evola. Has anyone lately tried to revive that idea? In particular, has anyone tried to harmonize it with the latest in anthropological, archeological, and historic research?

As for the second thesis, it answers the question why would the Nordics so readily adopt Christianity in the Medieval era? Perhaps they did recognize in it, or read into it, the echoes of a more ancient tradition. Gornahoor has supplied enough quotes to make that plausible. The competing answers, (1) the entire European continent was put under the spell of a global multi-generational Semitic plot, and (2) the ``Christians'' haulocausted the ``pagans'' or converted them by force. Option (1) is an insult to our ancestors while option (2) is logically flawed, since the Christians forcing the conversions were themselves Nordic. Evola wants to split the baby and extract the Ghibelline Middle ages from the Guelphs. But conceding so much to the non-Tradition seems unwise, when what is called for is another long march through the institutions. The full translation follows.

\begin{quotationx}
Wirth claimed he could reconstruct not only the history of the Nordic-Atlantic race, but also its religion. It would have been a higher, monotheistic religion, quite distinct from the animism and demonism of the black or Finish-Asiatic aborigines, without dogmas, of a great purity and potentially universal. At its base there would have been a type of natural revelation, that is, a perception of spiritual laws directly suggested by nature. When the Arctic freeze occurred, winter was prolonged for six months, so that the annual return of the sun had to be seen by those people almost as a liberation, as a resurrection of life. This is precisely the point of the winter solstice; the solar light appeared as a divine manifestation and the bearer of a new light; the year was the theatre of this manifestation and the winter solstice -- being the lowest point of the ecliptic, in which the light seems to permanently die, sinking into the earth or the waters, but instead miraculously rising from there -- was the decisive point of this cosmico-religious experience.

As we said, the sacred series for Wirth would have precisely fixed in the Nordic-Atlantic civilization the various phases of this symbolic annual event, summarized, in general, by the circle with a cross inscribed. The primordial religion of 15,000 BC would have therefore been solar and permeated by the sense of a universal law of eternal return, of death and rebirth. Like the light, so also the life of men has its ``year'', its perennial dying and rebirthing. The Christmas of the Christians, the birth of the Saviour at a date that fell in the period in which all the people celebrated the winter solstice, for Wirth, would be a distant fragmentary echo of this prehistoric religion. In general, Christianity would have originated from the tradition preserved among an Atlantic group of Galilee, a country rich with traces of the megalithic solar tradition. The most salient events of the life of Jesus, up to his crucifixion, that recovered the theme of the god-year, giver of life, nailed to the cross of the year, would be pure symbols of the Nordic-Atlantic tradition. So, Wirth speaks of a primordial Nordic monotheism and of a ``cosmic Nordic Christianity'' that would therefore date back to thousands of years before Christ, anticipating thereby Protestantism (which for him would only contribute to a re-Nordicization of that tradition) and would have had naturally nothing to do with the Jews.

The connection with ideas already entertained by Chamberlain and Woltmann is obviously established here, and, in addition, it has an imaginary point between a resumed tradition of early prehistory and the themes of dying and rising again and of eternal renewal so dear to German romanticism and the modern Faustian religion of life. Nevertheless, as to this last consideration, a divergence of views between Wirth and other racialist such as, for example, Gunther is quite visible. The concept of ``dying and rising again'', which for Wirth would make the keystone of the Nordic religion, Gunther would probably be able to carry a Semitic-Levantine spirit; and a divergence not less sensible remains in the fact, that while Wirth claimed that the symbol of a priestess or divine mother was at the first level among the Nordic-Atlantic people, who would even have called their land the ``Land of the Mother'', mo-uru, Gunther and various others related more sensibly similar conceptions to the meridional races and, at most, to the Celts, who would be a race already far from that pure Nordic and more akin to the Mediterranean races.

Besides, it is necessary to clearly distinguish the value and the significance of the Arctic thesis (or, as we prefer to call it, the Hyperborean) in itself from Wirth's arbitrary personal adaptations, because the plane to which it belongs is quite distinct and has a totally different dignity than these reconstruction of contemporary researchers, reconstructions, nevertheless, not lacking in interest as indications and obscure presentiments of a truth.

Von Leers writes that the preceding epoch of liberalism and scientism was characterized by three fundamental ideas.

\begin{enumerate}


\item The equality of the human race

\item Nordic barbarity and the origin of every civilization from the East

\item Finally, the Hebrew origin of monotheism
\end{enumerate}

These three ideas in the cycle that leads up to Wirth are defeated and overturned:

\begin{enumerate}


\item Humanity is differentiated into quite distinct races

\item Civilization did not come from the East, but from the North

\item It was not the Hebrews, but the Nordics who would have known, long before, a higher monotheistic religion.
\end{enumerate}
\end{quotationx}



\flrightit{Posted on 2011-10-22 by Aeneas}

\begin{center}* * *\end{center}

\begin{footnotesize}
\begin{sffamily}
\texttt{Charlotte Cowell on 2011-10-23 at 07:47 said:}

We might find out when the ice-caps melt. With respect to the Hebrews, Nordics etc, do you agree with Steiner that there are a number of different root races? Mainstream science reckons that everyone on Earth derives from (4?) African women, while DNA research has pinpointed around 223 human genes for which there's no link with our apparent predecessors in the animal world.

I still think the emphasis on Nordics is basically an inkling towards extra-terrestrials, if only at least in the sense of tapping into another dimension. That in itself is not unreasonable given that there are at least 11 known dimensions to our existence and cross-over is inevitable. Shambhala is in a different dimension to this one we in habit so it makes sense, and I also think it's reasonable to assume there are portals between them.

Sooner or later we should get onto Stargates, even if `they' (the ones who shouldn't be in charge) are busy trying to close them all up. Come on boys, form a circle around the stargate so we can keep on using it!!

\hfill

\texttt{EXIT on 2011-10-23 at 09:49 said:} 

I'm sure in your own mind, Exit, you think you have made a decisive point. Little, if any, of what you wrote was even mentioned, never mind, defended here. There is little that can be said to a mind that have been destroyed by propaganda. 

One can't rule out options 1 and 2 because they sound ``offensive'' to you. On the contrary, Christianity--which is spiritual slavery--is offensive to the ancestors of the national religions. This is what I said from the very beginning until you guys stole my quote and flipped it around. Gornahoor can't even come up with an original argument. One cannot say that Christianity is a true religion because it has survived this long. If that were the case then modernism and democracy are true as well. Conspiracy and evil, like it or not, is a part of humanity.

\hfill

\texttt{EXIT on 2011-10-23 at 10:06 said:}

\url{http://freethoughtnation.com/forums/viewtopic.php?t=2128}


As for those who don't believe Christianity was imposed on Europe click on the link, for starters. 

The question asked, Exit, even assuming your ``belief'' has merit, is who exactly did the ``imposing''? And who was the first imposer? There is a logical regress that manages to elude your mind.

BTW, whatever happened to the old Guenon/Matgioi saying, ``Love religion, hate the religions''?

\hfill

\texttt{Izak on 2011-10-23 at 13:23 said:}

``Gornahoor can't even come up with an original argument.''

I'm speechless.

\hfill

\texttt{HOO on 2011-10-23 at 16:13 said:} Charlotte Cowell ''emphasis on Nordics is basically an inkling''

An inkling?! The only thing rational, to re-trace, and re-manifest the species of Ultima Thule, is an inkling?

in·kling (ngklng)

n.

1. A slight hint or indication.

2. A slight understanding or vague idea or notion.

You are obviously (trying to) create an alibi for something you lack. Do better.

Unless you unconsciously express an obscure presentiment, relating to the etymology of an ``inkling'' from Old English ``nikken: to mark a text for correction''.

haha Charlotte Cowell ''the emphasis on Nordics is basically an inkling towards extra-terrestrials''

I'm not a UFO (yet), but I've sure felt like someone thrown onto earth (from what is extra to terra) since I was a child. Charlotte Cowell ''Come on boys, form a circle around the stargate''

And what is the stargate? Are you it? Izak: ''I'm speechless.''

Me too! Almost. Only thing I could think of is I wish I could think of something to heal the poor guy of misdirected potency.

He's never answered what he's trying to achieve.

\hfill

\texttt{Perennial on 2011-10-24 at 01:40 said:}

HOO, could you elaborate on your last comment? I am not clear as to whom or what you are refering.

\hfill

\texttt{Perennial on 2011-10-24 at 02:43 said:}

Not this again EXIT. I again find it difficult to give merit to such material when in contains-count em'- 0 references, at least none that I find. What is the allergy to referencing stuff?

\hfill

\texttt{Perennial on 2011-10-24 at 02:45 said:}

Forgive me, one book is mentioned. But still no citations.

\hfill

\texttt{Charlotte Cowell on 2011-10-24 at 05:47 said:}

Funny, after I wrote `inkling' I thought someone might accuse me of being a groupie of Oxford's famously literary mystics :-)

I don't think I need any kind of alibi and have ambivalent feelings towards the whole Nordic question that obsesses people here. It's just that having looked into various testimonies and ideas regarding Nordics it occurs to me that it's most logical to assume they're somehow extra-terrestrial. Could be wrong, could be right, who knows?

As for the stargate, I wouldn't describe it as a person but why not go through it (one of them) and see where it leads... . :-)

\hfill

\texttt{Charlotte Cowell on 2011-10-24 at 05:52 said:}

Especially if you're SERIOUSLY looking for Ultima Thule, that place beyond the borders of the known world... ..

\hfill

\texttt{Boreas on 2011-10-25 at 15:43 said:}

Charlotte: with the risk of sounding like a loose cannon, I firmly believe that Shamballa and the palace of Sanatkumara exists in the Gobi desert, in the etheric levels of the earth, and that the ``extra-terrestrials'' are the ones who have handed down the Primordial Tradition throughout the ages, beginning in ancient Lemuria. Once, when I was in a deep-seated spiritual turmoil, I had a glimpse of this place on the borderlines of my consciousness; very close to a dream state, but in full waking consciousness. It was a frightening and terrifying place for a ``normal'' man to come in touch with, with an aura of supernal light around it. No evil thought could live for a second in close proximity of it. I met one of its messengers in the desert steppes around the supernal temple. He had the characteristics of a hindu noble, a Kshatriya mounted on a horse. Everywhere was windy, and there was like a sound of a thunderstorm blewing through the desert. I dare not even imagine what secrets the palace itself holds.

I know, this sounds like second-rate fantasy, but I'm not kidding. Namaste.

\hfill

\texttt{Boreas on 2011-10-25 at 17:10 said:}

And as an interesting sidenote to my last post, I learned about the different theories of Shambhala (or Shamballa) and Sanatkumara the following day, when via a ``co-incidence'' I browsed through these in the web. The experience came before the theory, which was somewhat of a proof for myself personally at least. After this I learned that it was the time of the Wesak festival in May. Interesting to say the least. One things for sure = There's no smoke without Fire.

\hfill

\texttt{Charlotte Cowell on 2011-10-25 at 18:05 said:}

Boreas, I've also heard the Gobi desert theory (at least twice), so I don't think you're a loose cannon to say that. The extra terrestrial ideas are things I'm just starting to look into in earnest but there are some extremely credible testimonies with respect to `life from other planets/dimensions' in relation to all aspects of our existence, and one would have to be stubbornly and willfully blinkered in order to refuse to explore these. (The Australian girl Tracey Taylor is one I heard recently and her experiences with Horus -- who He/she equates with Quetzalcoatl -- are rather mind boggling.

Going back again to `location' my own experience involved seeing a very distinct mountain range that I took to be the Himalayas, although I periodically surf the web for photographs of icy mountains to see if I got that wrong. The truth is I'm not sure, although I found two locations that seemed to be the best possible fit from a visual perspective. One was in the Himalayas and the other was on the border with China -- I've read about this specific place twice and forgotten it twice, only to have to hunt for it all over again! I've also read an enthralling account of someone who actually found Shambhala in the Gobi desert and described an utterly extraordinary experience -- possibly it's the one you read? Again, I'm darned if I can remember, how annoying of me! I shall try to revise this and see if I can find it, possibly it was one of the books on Agartha. Your experience sounds similar to the one I read about in many respects and I'm with on your method of personal confirmation -- I also have experience first and then find it backed up later in writing or from other accounts. It's a failsafe way for the the self at least, although doesn't much help in convincing others!

By the way, my own Shambhala experience also coincided with a phase of profound spiritual crisis -- personal crisis on all levels actually... .no pain no gain as they say. I'd love to know more about the different theories you gained insight into if you've got time to relate them here (or refer me to somewhere else they may be written).

The people I saw in the place I went to -- gathered around a camp fire -- did not resemble Brahmin type character you mention. They looked more obviously Tibetan and were dressed in quite ordinary nomadic clothes. They were gathered around a campfire that I actually fell onto from a great height with an almighty crash. It's a long story all this so I won't go into it, except to say that eventually one came forward and offered to guide me -- I got the impression they'd decided someone had to seeing as I was there and that nobody was too pleased about it.

I do, however, know the wind you describe. I've also experienced the intense `psychic wind', as I call it, and on each occasion it has brought an almost unimaginable terror that would have driven me into blind insane panic if I'd been in rather than out of the body. As you are no doubt aware, emotions don't register in the same way in that state due to separation from the body and the accompanying intrinsic knowledge of immortality. Yet still one knows that one really should be mortally afraid.

Tangentially, it has occurred to me in passing that this wind may be connected in some way with the Earth's passage -- at very rapid speed after all -- through space, of which we are otherwise blissfully unaware. But also high adepts will use it to frighten people in order to test them.

As you say, no smoke without fire. Cx

\hfill

\texttt{Charlotte Cowell on 2011-10-25 at 18:14 said:}

PS, regarding the keepers of the primordial tradition, it's perhaps worth bearing in mind that many contactees, credible or otherwise, describe there being a `council'. Ancient India teaches us of the nine wise men, perhaps they are analogous to this council?? Horus continually resurfaces in this context, ditto for Quetzalcoatl, although Tracey was the first one I heard describing them as one and the same. When she emerged from the other dimensional experience she was left with the profound phrase `Architect of Consciousness', and from that moment on began to draw incredible, mathematically brilliant drawings and produced reams of text in some unknown language resembling Sumerian. But what impressed me most was the beautifully serene spiritual light she radiated... .more and more people will be born that have that same light in abundance, indigo children are just the start.

\hfill

\texttt{Charlotte Cowell on 2011-10-25 at 18:43 said:}

Boreas, I just found this, which is brief but seems to tie in very well with what you described. Also, the 8 petal-lotus palace and the 8 different regions perhaps explains why you and I saw things from a different perspective? Also, I'm drawn to comment on what Alexandra David-Neel says about the individual she saw moving at lightening speed and looking up... .I arrived from above (on the back of an eagle). The nomad-type people moved normally until we got inside the mountain, whereupon they did indeed zoom along at astonishing speed. At the time I assumed it was so I couldn't get a good look at what was inside, although I had a definite impression that it was something utterly extraordinary and very secret. There was a lot of light in there. The central hall was awesome and filled with seemingly countless beings, all holding lamps, queuing upwards in a long snaking line for an audience with some tremendous being hidden from view. I could not see the faces of any of the beings, who were all cloaked and holding lamps -- at the time I wondered if this was because they were dead, as in fact I was quite beyond my knowledge zone and passively observing without really understanding what was happening. The guide told me I had to be silent but said it was a meeting of saints and angels. I had no clue of who or what they were meeting with and nor was i subsequently able to figure this out.

\url{http://www.dreamscape.com/morgana/juliet2.htm}


This also looks like a good resource, I just found:

\url{http://www.kalacakra.org/}


Cx

\hfill

\texttt{Boreas on 2011-10-27 at 06:30 said:}

Thanks for your thoughts and links, Charlotte. My own literary reference was, if I remember correctly (it's been years since the occurence happened), Guénon's `King of the World', which I had read at the time and had come up with the name of Sanatkumara. I then remembered that Evola had also dealt the concept of a chakravartin somewhere in `Revolt... ', in relation to Melckhizedek and Sanatkumara. After this I studied the matter simply from internet, wikipedia etc. and thereon stumpled into different points of view about the matter, everything from traditional hindu and buddhist points of view to theosophical and neo-theosophical points of view. Despite many confusions and distortions caused by some new age authors who cannot possibly be taken seriously, all these sources nevertheless seem to point into the same direction. In these matters one most certainly has to ``test the spirits'', so to speak, in order to get through all the ``smoke''.

\hfill

\texttt{Charlotte Cowell on 2011-10-27 at 12:26 said:}

Thank you for clarifying that Boreas -- I've not read those texts you mention but will look into them now. Fascinated to see how Melchizedek might tie in with Sanatkumara... .I must admit I've tended to give scant credence to stories of Sanatkumar, perhaps because of unjustifiable prejudice against HPB. She poses a LOT of difficulties with respect to discernment, but then so does Gurdjieff (I steer clear of both although they are fascinating personalities and evidently pioneering spirits). Incidentally Madame Blavatsky's Baboon is a great and very funny book for the down-to-earth view of these characters.

My own experience happened in 2009, when I was engaged in serious meditative work and had been for almost a year by that point, using the Rosicrucian method. Someone I'd been corresponding with introduced me to the work of Roerich, which I immediately fell in love with and -- tending not to do things by halves -- spent several hours that very day pouring over his paintings.

However, when I did not `set out' with the intention of visiting Shambhala, although clearly it must have been in my subconscious somewhere because of all the references in Roerich's work. (I had a take-it-or-leave-it attitude to such mysteries, which by the way is probably helpful when it comes to gaining access).

I was actually in search of the Midnight Sun, notions of which had been absorbing me for some time, partly because of Mark Hedsel's extraordinary work, The Zelator, a very highly charged text as you will no doubt be aware if you've read it. This was an exercise that would probably have been quite frightening had I not been foolhardy (I have learned better since then, having ventured into terrain I most definitely did not want to be in, attracting all kinds of terrible things by my rash behaviour, but I'm glad I managed to experience certain things before coming to my full senses, and fearlessness has its own advantages for sure).

In any event, the entire procedure that ended up in Shambhala was very complex -- involving splitting of consciousness and invocation and protection of the angel Gabriel (with satan right behind, I will admit) -- followed by reunification of consciousness with the angelic force. This was followed by a long ascension upwards, bearing in mind we already were at cloud level. How I (we at this stage) got to this point from the interior sun I don't know, because I blacked out when I got near to this and still don't know what happened in that in between phase, I think it was rather frightening. In any event it all but finished off part of me, which was put back together by a small group of beings on the cloud who wrapped that part in a rainbow cocoon, from which emerged the winged butterfly consciousness.

The two parts of myself (object and subject) then came together and the butterfly aspect began transforming into different white birds as we floated upwards, culminating in the eagle, whose back I rode on. (This was usual for me, I frequently would travel with an eagle and recommend this as a very safe and accurate means of transport).

anyway it was freezing cold and we flew for a long way until at a certain moment (our heads were very close together, I was crouched over it) I wondered where we were headed. That's when I saw the large triangular ice-capped mountain appear and was given to understand it was the himalaya. I still had no sense of why we were going that way, although by this point I had somehow figured out that we were looking for someone -- rather than something.

We flew over the mountain (it was spectacular) and then I became aware we were earnestly looking for some kind of community or outpost. I must have asked some question about how we were meant to see them because the eagle at that point went much lower. By the way it was a crystal clear blue-sky (usual for this state), very bright, we could see for miles.

Very soon after dropping lower we saw a large campfire, which stood out like a sore thumb against the snow. This was fairly exciting to see a sign of life (I still didn't consider whose) but both myself and the eagle had a moment of questioning; was it safe to go visit them, would they be friendly and -- if I did want to see them -- what was the best way of getting down there?

Having got so far the die was cast, I was going down there -- but how? A split second later either the eagle or myself and the eagle swooped down at lightening speed, plummeted in fact, causing me to wonder later if the eagle had in fact been shot down. At the time I felt as if I'd been tipped off its back.

I landed with an almighty crash right on top of the extremely large camp fire and was immediately disconcerted (to say the least) to discover that all was in darkness apart from this fire; it was in fact the dead of night. I had a maelstrom of thoughts and even feelings -- in so far as one can have feelings in that state -- worrying about the eagle, fearful of the dark, wondering about the fact I was in a fire, and gradually becoming aware of being silently watched.

The campfire was surrounded by people, all men I think, though I could be wrong about that, sat quite some way back, too far for me to see them properly, although their eyes glinted in the darkness. As my eyes adjusted I saw more of their appearance but that wasn't my main concern. Mainly I was concerned about being really worried and stuck in the middle of nowhere with all these men who seemed sinister and unfriendly.

After sometime of being totally mesmerised by the fire and looking up wondering about the eagle I began to make a tremendous fuss. I wanted to escape and at lightening speed ran as far away as I possibly could from the fire and its strange people. In no time at all I was at the edge of the camp, but the only thing that really marked it out as being the edge was the fact that the circle of light from the fire only reached to a certain diameter. Having reached this point and seeing that beyond was pitch-black, I found myself in a catch 22 situation. In or out? On balance out seemed scarier due to the absence of light (Yetis??) and so I sat right on the edge and basically sulked for quite a while.

Sooner or later I got bored and went back. I'll cut the story in a sec as it's becoming massive but after a short while one of the men stepped forward -- I sensed very reluctantly indeed -- and basically said he would be my guide. I got the feeling someone had to be seeing as I was there. I'm not sure they liked the fact I was female, In any event, he set me an exercise to do with fire and snow (involving diving in and out of fire and snow drifts and whirling round as quickly as possible) that lasted for the rest of the entire night, absorbing me completely in quite obsessive, pretty much lunatic-seeming fashion, until the fire died down and and a grey dawn broke. (this was apparently something to do with the creation of stars).I finished my fire exercise by making the formerly white/yellow roaring flames make a relatively gentle orange `Shin' sign with me, then stepped out of it.

Dawn revealed a more mundane scene. The people looked like ordinary nomads and were dressed as such. They quickly made preparations to leave the fire place and began walking rapidly towards the mountain that I now saw was behind us. It felt like maybe 500 metres away and as we went closer (the guide was walking with me at the back of the line) I saw there was a large crack and that we were clearly going to walk inside. I was apprehensive as I don't like enclosed spaces (or the dark, not at all, I'm actually quite afraid of them in real life) and didn't know where we were going. The guide told me not to worry. That's when we went inside and began whizzing through.

\hfill

\texttt{Charlotte Cowell on 2011-10-27 at 12:33 said:}

`ordinary' life I should say, especially enclosed spaces, although for some unfathomable reason I dreamed about being shut up inside a bread bin the other night and there was actually another person in there with me!! It was so cramped I lifted the lid from the inside and climbed out, waking up to find myself scrunched up into a tiny ball, every muscle in my body aching like mad :-)

\hfill

\texttt{Charlotte Cowell on 2011-10-27 at 12:42 said:}

It was my own bread bin by the way, the one that's in the kitchen

\hfill

\texttt{Charlotte Cowell on 2011-10-27 at 12:45 said:}

Funny, I've never read a thing by Guenon (sorry, I know I should as he is the main focus of this blog), but having quickly done a Google search we seem to have highly convergent areas of interest, I'd probably really enjoy some of his writing... .

\hfill

\texttt{Max on 2015-03-11 at 17:39 said:}

Regarding archeological evidence of the sacral architecture of the Nordic tradition, for the starting point of a study could be used the gothic cathedrals of the middle ages that almost ``appeared out of nowhere''. Of course they did not, but the sources probably lie in earlier wooden construction or even sacred groves of living trees. To get some ideas of what that was like the look up Norwegian stave churches. There are also still wooden churches or temples in Russia, India and the Himalayan region. The next question is why wood was replaced with stone, obviously it represents the combination of the Roman and Germanic traditions, but maybe there are deeper reasons as well. A breathing grove of broad leaf trees will also take on another meaning in a land otherwise dominated by dark coniferous forests. Sacred groves does not make much sense in the jungle.

\hfill

\texttt{Marcel K. on 2015-03-11 at 21:14 said:}

The arctic myth is hard to defend because all existing theories regarding the Proto-Indo-European homeland point to the east, with the most northerly possible possible origin being in the Pontic-Caspian steppe in southern Russia and central Asia. It was picked up by the Günther-Woltmann-Chamberlain school of anthropology, which had great influence in the 20th century, and in all probability lead many thinkers to interpret the arctic myth all too literally.

As I pointed out earlier, even Evola adopted the belief that the European tradition was spread by blonde, blue-eyed dolichocephals who invaded inferior races from the north and gradually assimilated into existing, non-European populations which resulted in racially mixed populations of today.

This is false because pre-European farmer and hunter-gatherer populations exhibited typically European traits thousands of years before PIE speakers had them. The earliest human who exhibited both blond hair and blue eyes was a pre-Indo-European farmer from Hungary who lived 6,000 B.C. . Long before that the Hunter-Gatherers native to Europe spread the genes for blue eyes, while light skin arrived with Middle and Near Eastern farmers from the south. In contrast the people who may have brought PIE into Europe 4,000 B.C. would have been of a rather dark `Eurasian' Caucasoid type, or a near eastern type depending on the theory one subscribes to.

2015 has been a very exciting year in anthropology thus far and with the new paper by Haak et al. the `steppe hypothesis' gained a lot of support. In this paper the researchers posit that the Yamna culture native to the Pontic-Caspian steppe migrated into Europe and spread PIE. Quite a few modern populations derive a large percentage of their genes from the Yamnaya steppe people according the the researchers. What makes the steppe hypothesis so alluring is the fact that the archeology bears it out: the inhabitants of the steppe domesticated horses, developed the chariot and formed warrior bands with a Männerbund-like character. In short, they would have been perfectly capable of overwhelming the sedentary farmers and thinly-spread out Hunter-Gatherer populations native to Europe. Even in historic times steppe peoples posed a significant threat to more sedentary populations.

In a hypothetical homeland in the far north however the only possible mode of subsistence would have been that of the native European Hunter-Gatherers in Scandinavia and Eastern Europe. They would have been spread out very thinly and their cultural influence would have been extremely limited. Neither would they have been capable of the military organisation required to spread PIE far and wide by elite dominance. If a Hunter-Gatherer culture of the north would have transitioned into a more advanced culture capable of dominating populations all over Eurasia, they would have left an archeological record.

Of course if the Yamnaya people were the Proto-Indo-Europeans one could hypothesize a more northernly homeland for the cultures that preceded them. However, the Samara culture of 6,000 B.C., the predecessor of Yamna, was situated in the south of Russia as well. If there existed a culture further in the north which preceded Samara it has left no traces. The probability for a southern source is much greater in this case.

An interesting event that roughly coincides with the early developments of PIE is the Holocene climactic optimum. If the early Indo-Europeans to whom the arctic myth would be applicable lived somewhere near Siberia, for example, where a warming of up to 9°C occurred in a relatively short time, they would have experienced a transition from an environment that was arctic in the metaphorical, not in the literal sense, to a drastically different environment even without migrating. Perhaps this extreme change in climate could account for some of the descriptions in the Vedas.

It would be interesting to investigate links like this and match the descriptions found in traditional sources against the wealth of scientific evidence that resulted from the advances in anthropology and archeology in the last two decades.

\hfill

\texttt{Aeneas on 2015-03-12 at 00:06 said:}

Of course, Marcel, Evola rejects out of hand any consideration of profane science or archeology. To adhere to science would be considered a restriction to free thought. Moreover, profane science can never reach into the spiritual and soul qualities of any ancient groups it studies. Whether that is a position that can, or should, be maintained is another issue.

\hfill

\texttt{Cologero on 2015-04-19 at 09:57 said:}

Marcel, we keep coming back to this question: is positive science the ultimate source of all knowledge? If ``yes'', then this whole exercise is pointless. So anthropologists claim to go back 8,000 years based on DNA of a few dozen samples. All we can say is that the trail grows cold beyond that point. As an mental exercise, consider a Swedish woman who went to Jamaica where she was impregnated by a native, and gave birth to a daughter. For generations after, her mitochondrial DNA would leave a trace in the Jamaican population, misleading future anthropologists.

Evola's view, and really that of anyone with a spiritual view of life, rejects the totalitarian nature of science. Science presumes a physical cause for everything, a universal human nature, and so on. Suppose, instead, that is not true. Then no DNA evidence or linguistic analysis can reveal the soul and spirit of any human group. Differences in consciousness are invisible to science, and so on and on. But we are just repeating ourselves.

Scientists reject out of hand a ``spiritual'' anthropology: that is, the analysis of myths, rites, revealed, texts, symbols, and so on, that reveal the inner nature of the population. It is necessary to make a choice.

\hfill

\texttt{White Rabbit on 2015-04-20 at 03:20 said:}

sweden yes

There is actually an interesting passage from Guénon which relates to this question.

``This correspondence between higher and lower holds good for historical facts no less than for anything else; they likewise conform to the law of correspondence just mentioned, and, thereby, in their own mode, translate higher realities, of which they are, so to speak, a human expression. We would add that from our point of view (which obviously is quite different from that of the profane historians), it is this that gives to these facts the greater part of their significance.

The symbolical character, while common to all historical events, is bound to be particularly clear-cut in the case of of events connected with what may be called `sacred history'; thus it is recognizable in a most striking way in all the circumstances of the life of Christ.

If the foregoing has been properly grasped, it will at once be apparent not only that there is no reason for denying the reality of these events and treating them as mere myths, but on the contrary that these events had to be such as they were, and could not have been otherwise; it is clearly impossible to attribute a sacred character to something devoid of all transcendent significance.''

(Symbolism of the Cross, Preface)

\hfill

\texttt{( ) on 2015-05-01 at 05:57 said:}

Why the preoccupation against a Semitic origin of monotheism, as we find in the scriptures? History seems to show monotheism born from the crucible of the Babylonian Captivity, where ``YHWH is the highest of the gods'' was realized to be ``YHWH is the only God,'' as the quite easy conquering of the Jewish people seemed to make the Babylonian gods look vastly superior. One need only look at the development of Isaiah's later parts in the context of the Old Testament to see this.

That's not to say there wasn't perhaps a Hyperborean monotheism- after all, linguistics seems to show a ``deus pater'' who the proto-indo-europeans worshipped. But is the semitic development illegitimate?

I notice comments bellow that suppose ``Atlantic Galileans,'' which seems to me a quite ludicrous assumption. Christ was semitic- in a Hellenized culture- and quite proudly so, as was Paul. The Romans did not seem concerned with the Jews and messianic Jews (later Christians) from a racial perspective; though as we know the early Christians spent quite a bit of time trying to figure out how to spread to other races.

It seems bizarre to me to read some Europeanness into the thoroughly semitic origins of Christianity. That Christianity found its traditional form and height of beauty among each of the great civilizations, the Romans, Byzantines, and even the less grandiose but equally traditional Armenians, Coptics, Ethiopians, and such, is obviously true. But is there really a need to wipe the Semite presence out of our real- and quite Traditional- history? Is this antisemitism and its accompanying baseless paranoia at its worst?

If I'm misreading and, in fact, nobody is doing this, I'm glad and apologize for my misinterpretation and waste of your time.

\hfill

\texttt{Cologero on 2015-05-01 at 07:20 said:}

Mr Open Brackets, thank you for your thoughtful comments. Keep in mind that every dialog has a context and an intended audience. So that post was addressed to the neo-pagans. I'm afraid, however, that we are at odds about the historical method you employ. Specifically, ``monotheism'' cannot be ``born'', if you mean that a people simply adopted an arbitrary idea. As Guenon pointed out, ``monotheism'' is part of Tradition, not just that of the Semites. Please read his writings on that topic and tell us your objections to it.

Furthermore, we addressed the concept of ``Europeanness'' already, since there was no ``Europe'' to speak of. It itself was born from Christendom, and before that there were simply groups of different races inhabiting a land. Now if you are objecting to the mixing of Greek thought at the ``origins'' of Christianity, your argument needs to be directed to church fathers, medieval scholastics, and even recent popes. I'll not go into all those same arguments all over again.

Now if you ``see'' antisemitism and ``baseless paranoia at its worst'' in this, then you need not apologize for wasting my time. Rather you need to apologize for libel, a much more serious offense.

\hfill

\texttt{( ) on 2015-05-01 at 09:47 said:}

I apologize, as it's clear I should've been more direct about what concerned me, because I agree with all you have to say about the idea of ``Europe'' being borne fundamentally of Christianity, and I have no trouble with the incorporation of Greek thought into Christian thought, a process that was truly natural- I have read your writings on both of these issues and I have no criticisms of either.

My complaint about antisemitism applied only to (and I should have been clearer about this in the first comment; this truly was my mistake and I apologize) Evola's claim ``Christianity would have originated from the tradition preserved among an Atlantic group of Galilee,'' which is entirely unsubstantiated and unnecessary. What is the point of conjuring up some Northern group in the Holy Land, if not to baselessly deny semitic influence on Christianity and Christ? One need only read the Bible to see how much Judaism and its history inhabit the roots of the Christian tradition.

As to the question of monotheism in Judaism, I do not think that monotheism was simply ``born'' and that was not my intended implication. Rather, I think it is clear that the Jews were at one point pagans and that they became monotheists -- much as how Pagan Europe was Christianized naturally -- at a specific moment in history. My point in even bringing this up was again to contest that there need be an ``Atlantic'' presence for the Jews to find Tradition when they obviously did on their own. If you disagree with this understanding of things, I can point to the Old Testament passages where the Jews are clearly pagans (albeit strict in their henotheism) who understand YHWH to be their national deity among many, and then to the point in Isaiah where the insistence changes. But this point is not my primary concern, which can be found above.

\hfill

\texttt{Cologero on 2015-05-01 at 12:57 said:}

Thanks for the clarification Mr ( ). Evola probably wrote too much and I don't even recall that passage. Recently we offered this translation:

\begin{quote}\footnotesize

In its original sacred form, prior to the period of the prophets (that indicated the first mystical and democratic fall of Israel's ancient tradition), the idea of the Messiah had many traits in common with familiar conceptions and ideals of essentially Aryan civilizations, from which, moreover, the Hebrews more than once borrowed many elements.
\end{quote}


That certainly is disputable, but perhaps the influences really went in both directions. But my intent in publishing that was not to promote anti-semitism but rather to dispute it. If the Hebraic idea of the Messiah was original, then it resonated with a similar ideal in the West.

I am not really interested in your OT ``proofs'' since I don't accept Sola Scriptura as a doctrine. Sacred writings need to be read on multiple levels.

\hfill

\texttt{( ) on 2015-05-01 at 14:24 said:}

If Evola's point there isn't greatly important, I'm not worried about it, then. I may've been too concerned about a minor issue that struck me as odd, here, and in that case an apology for wasting your time would, in fact, be appropriate.

I also do not accept Sola Scriptura and was not trying to use the sort of cheap scriptural ``proof'' that one sees frequently in highly-modern Protestant discourse; you'll notice I never actually did offer a proof, just to show you what I was referring to. I was merely speaking about the scriptures inasmuch as they are historical documents that give us insight into the time and place where they were composed- I'm sure you'd agree that this, and the work of the historian as a whole, is not useless. Obviously the sacred writings are useful and truthful on many more levels than this, and should never be restricted to a single reading- as the Rabbinic adage goes, there are ``seventy faces to the Torah;'' likewise with all scriptures.

But again, the matter I was referring to is not pressing, and I'm not trying to press it. My concern was for Evola's strange claim and you have addressed it.


\hfill

\end{sffamily}
\end{footnotesize}

